% =====================================================================
%  5 -- STIFF INTEGRATION OF THE AUGMENTED FLUX STATE  (source node S10).
%
%  Source: tier2.md Part IV in full and in source order -- header
%          (3660-3664), IV.1 (3668-3699), IV.2 (3703-3747),
%          IV.3 (3751-3776), IV.4 (3780-3811), IV.4b (3815-3867),
%          IV.5 (3871-3938), IV.6 (3942-3966), IV.7 (3970-3994),
%          IV.8 (3998-4060), IV.8b (4064-4205), IV.9 (4209-4251),
%          IV.10 (4255-4285), IV.11 self-check (4289-4306).
%  Nothing is re-homed into or out of this range: the source's own
%  order already runs stiffness -> method -> state -> identifiability
%  -> Jacobian -> conservation -> error control -> energy -> cost.
% =====================================================================

\section{Stiff integration of the augmented flux state}
\label{part:integration}

\begin{repopointer}
\textbf{Files:} \code{fluxes/integrate.py} (265 lines). \quad
\textbf{Oracles:} \code{tests/test\_integrate.py} ---
\S\ref{sec:codes10}. \quad
\textbf{Notebook:} \code{11-integrator-and-handoff.ipynb}. \quad
\textbf{ADR:} 0006 (reference integrator).
\end{repopointer}

The last node of the tier, and the one that produces something the shipped
dataset does not contain: \textbf{per-reaction integrated fluxes \Phifl}, the
labels Target A needs.

% =====================================================================
\subsection{Stiffness, from the Jacobian}
\label{sec:stiffness}

The network's Jacobian is a product of the two objects this tier is about:
%
\begin{equation}
J_{ij} \;=\; \frac{\partial \dot Y_i}{\partial Y_j}
\;=\; \sum_k \nu_{ik}\frac{\partial R_k}{\partial Y_j}
\;=\; (\nu D)_{ij}, \qquad
D = \frac{\partial R}{\partial Y}\in\mathbb{R}^{r\times n}.
\label{eq:jacobian}
\end{equation}
%
Its eigenvalues are (minus) inverse relaxation timescales. The spectrum here is
spectacularly spread, and the spread is \emph{measured} rather than estimated:

\begin{itemize}
\item At the shortest label step $\dd t_1 \approx 1.011\times10^{-6}$\,s,
  \textbf{99.99\%\,/\,100.00\%} of (state, isotope) cells on the training grid
  have $\tau_i = Y_i/|\dot Y_i| < \dd t_1$ \resultsref{2026-07-10}. So the
  fastest modes are $\ll 10^{-6}$\,s.
\item Trajectories are integrated out to ${\sim}10^{8}$\,s.
\end{itemize}

That is a \textbf{stiffness ratio of at least $10^{14}$} within a single problem
--- and it is not an artefact of a pathological state, it is the generic case
here. Tier~0 \S0.7's timescale hierarchy, now with a number attached.

\paragraph{Why this kills explicit methods.} Forward Euler on $\dot y = \lambda
y$ is stable only for $|1 + h\lambda| \le 1$, i.e. $h \le 2/|\lambda|$ for real
negative $\lambda$. With $|\lambda| \sim 10^{6}$\,s$^{-1}$ the step is capped at
$2\times10^{-6}$\,s \textbf{regardless of how smooth the solution is}, and
reaching $10^{8}$\,s needs ${\sim}5\times10^{13}$ steps. The fast modes are long
dead --- they contribute nothing to the answer --- and they still dictate the
cost. That is the definition of stiffness: \emph{step size limited by stability
rather than by accuracy.}

And \S\ref{sec:slowmanifold} says why the fast modes die: they are the
relaxation onto the QSE manifold. \textbf{The stiffness and the equilibrium
structure are the same fact.}

% =====================================================================
\subsection{Why BDF, and what L-stability buys}
\label{sec:bdf}

An implicit method evaluates $f$ at the new point, so the linear stability
function does not blow up for large negative $h\lambda$. The relevant hierarchy:

\begin{itemize}
\item \textbf{A-stable}: stable for all $\mathrm{Re}(h\lambda) < 0$. Backward
  Euler and the trapezoidal rule are; \textbf{Dahlquist's second barrier}
  \citep{dahlquist1963} says no linear multistep method of order $> 2$ can be.
  So high-order multistep methods must settle for A($\alpha$)-stability.
\item \textbf{BDF-$k$} \citep{curtiss1952, gear1971}:
  $\sum_{m} \alpha_m y_{n-m} = h \beta f(t_n, y_n)$. Zero-stable for $k \le 6$;
  A-stable for $k \le 2$; A($\alpha$)-stable with shrinking $\alpha$ for
  $k = 3\ldots6$ ($\alpha \approx 86\degree$, $73\degree$, $52\degree$,
  $18\degree$; \citealp{hairer1996} --- the $k = 6$ wedge is so narrow it is
  rarely used). scipy's \code{BDF} is variable-order 1--5 with adaptive steps,
  which is the right compromise: it drops order where the $\alpha$-wedge bites.
\item \textbf{L-stable} \citep{ehle1969}: A-stable \emph{and} the amplification
  factor $\to 0$ as $h\lambda \to -\infty$. Backward Euler has
  $R(z) = 1/(1-z) \to 0$~\yes. The trapezoidal rule has $R(z) \to -1$~\no{} ---
  it is A-stable but leaves stiff transients \textbf{ringing} with undamped
  sign-alternating error.
\end{itemize}

\begin{invariant}
\textbf{L-stability is the property that matters here}, and it is worth being
explicit about why: the physical content of a stiff transient in this system is
\emph{relaxation onto the QSE manifold}. A method that keeps the transient
oscillating at amplitude 1 does not merely lose accuracy, it produces a
composition that oscillates around the manifold --- and since abundances near
the manifold can be many orders below the dominant species, the oscillation goes
negative.
\end{invariant}

Which is why \code{integrate.py} clips $Y$ at 0 inside rate evaluation:

\begin{srclisting}
Yc = np.maximum(np.asarray(Y, dtype=np.float64), 0.0)
\end{srclisting}

with the docstring's justification --- ``transient small negatives from the
linear multistep are not propagated into powers''. Note the clip is applied
\textbf{only inside rate evaluation}, not to the solver state: clipping the
state would inject mass and break conservation. Clipping the rate argument is a
regularisation of the RHS, and it is safe because the returned composition is
reconstructed from \Phifl{} (\S\ref{sec:solverconservation}) rather than from
the clipped values.

\textbf{Radau IIA} (implicit RK, order 5, L-stable) is available for
cross-checks, and the cross-check was run: BDF $\equiv$ Radau $\equiv$ BDF(rtol
$10^{-10}$) to \textbf{$5\times10^{-8}$} rel on the stiffest test state.
\resultsref{2026-07-12} Two structurally different integrators agreeing is the
only honest way to establish that neither is lying --- \foreignreg{1}{R-23}'s
Fortran-probe logic, applied to the ODE side.

% =====================================================================
\subsection{The augmented state \texorpdfstring{$[Y, \Phifl]$}{[Y, Phi]}}
\label{sec:augmented}

\code{integrate\_state} integrates
%
\begin{equation}
\frac{\dd Y}{\dd t} = \nu R(Y; T,\rho), \qquad
\frac{\dd\Phifl_j}{\dd t} = R_j(Y; T,\rho),
\label{eq:augmented}
\end{equation}
%
with $u = [Y, \Phifl] \in \mathbb{R}^{n+r}$ and $\Phifl(0) = 0$.

\Phifl{} is a pure quadrature: $R$ depends on $Y$ only, never on \Phifl. So the
\Phifl{} block adds no dynamics --- mathematically it is $n_{\rm rxn}$ running
integrals carried alongside. What it adds is:

\begin{itemize}
\item \textbf{exactness}: $\Delta Y = \nu\Delta\Phifl$ becomes an identity of
  the \emph{discretisation}, not just of the continuum equations, because both
  are advanced by the same solver using the same $R$ at the same stages;
\item \textbf{cost}: the state vector grows from $80 \to 687$ (\msn{80}) or
  $151 \to 1669$ (\msn{151}), and --- more consequentially --- the error test
  applies to all \Phifl{} components too (\S\ref{sec:errorcontrol},
  \S\ref{sec:costwall}).
\end{itemize}

\textbf{Gross, not net.} $\dd\Phifl_j/\dd t = R_j$ is the \emph{gross}
per-column rate, matching the engine's \fp{} convention exactly, so that the net
view is recovered downstream through \code{pair\_col} in precisely the same way
\code{store.py} does it (\S\ref{sec:store}). One convention, one place to be
wrong. The docstring is explicit that this mirrors the flux engine, and
\code{test\_rhs\_matches\_engine} pins it to $\le 10^{-12}$ rel.

% =====================================================================
\subsection{Why \texorpdfstring{\Phifl}{Phi} cannot be recovered from
            \texorpdfstring{$\Delta Y$}{dY}}
\label{sec:phinotrecoverable}

\begin{warn}
\warnglyph{}~This is the reason the node exists, and it is a one-line consequence
of \S\ref{sec:rank}.
\end{warn}

Given a target $\Delta Y$, the solutions of $\nu\Phifl = \Delta Y$ form an
affine set of dimension
%
\begin{equation}
r - \rank(\nu) \;=\; 607 - 79 \;=\; \mathbf{528}\quad(\msn{80}),
\qquad 1518 - 150 \;=\; \mathbf{1368}\quad(\msn{151}).
\label{eq:fibredim}
\end{equation}

So:

\begin{result}
\textbf{The shipped Zenodo labels (initial and final $X$, per state per
$\dd t$) do not determine \Phifl.} Not ``determine it noisily'' --- do not
determine it at all, up to a 528-dimensional subspace.
\end{result}

Consequences, each of which shows up somewhere else in the repo:

\begin{enumerate}
\item \textbf{Target A cannot be supervised from the shipped labels alone.} A
  loss on $\Delta Y = \nu\hat\Phifl$ constrains 79 of 607 directions. The
  remaining 528 are free, and a network is free to put arbitrary structure there
  --- including physically absurd flux patterns that happen to produce the right
  $\Delta Y$. Hence \code{integrate.py}.
\item \textbf{The minimum-norm pseudo-inverse $\Phifl = \nu^{+}\Delta Y$ is not
  the physical answer.} It is \emph{a} solution, the one of smallest Euclidean
  norm, and there is no reason the physics should choose it. (Near QSE the
  physical \Phifl{} has enormous nearly cancelling components --- precisely a
  \emph{large}-norm solution.) Anyone tempted by \code{np.linalg.lstsq} as a
  shortcut should reread \S\ref{sec:vacuous}'s $\kap \approx 1$ vs the relaxed
  continuum.
\item \textbf{Conversely, the redundancy is what makes masking meaningful}
  (\S\ref{sec:rank}) --- the same fact, read the other way.
\end{enumerate}

% ---------------------------------------------------------------------
\subsection{The flux head is unidentifiable under a
            \texorpdfstring{$\Delta Y$}{dY}-only loss}
\label{sec:unidentifiable}

\S\ref{sec:phinotrecoverable} drew the \emph{label} consequence of the 528-
(1368-) dimensional solution set: \Phifl{} is not recoverable from $\Delta Y$,
hence \code{integrate.py}. There is a second consequence, about
\textbf{training}, and it is arguably the sharper one.

Suppose Target A is trained with a loss on $\Delta Y$ alone (which is what the
shipped labels permit). Then the loss depends on the flux head's output only
through $\nu\hat\netflux$, so:
%
\begin{equation}
\nabla_{\theta}\mathcal{L} \;=\;
\Bigl(\tfrac{\partial \hat\netflux}{\partial\theta}\Bigr)^{\!\mathsf T}
\nu\Tr\,\nabla_{\nu\hat\netflux}\mathcal{L},
\label{eq:gradnull}
\end{equation}
%
and \textbf{every parameter direction that moves $\hat\netflux$ within
$\operatorname{null}(\nu)$ receives exactly zero gradient.} The optimisation
problem is rank-deficient by 528 dimensions out of 607.

What actually happens to those directions in practice:

\begin{itemize}
\item they are set by \textbf{initialisation} and moved only by weight decay,
  regularisation, and the implicit bias of the optimiser --- i.e. by everything
  except the physics;
\item under SGD they can \textbf{random-walk}, since noise in the gradient
  estimate has components there even when the true gradient does not;
\item and they are \textbf{invisible to every validation metric defined on
  $\Delta Y$}.
\end{itemize}

Three consequences:

\begin{enumerate}
\item \textbf{Target A's per-reaction interpretability is not guaranteed by a
  $\Delta Y$ loss.} The selling point --- ``a flux error is attributable to a
  named reaction'' (\S\ref{sec:bet}) --- requires the flux head to be predicting
  \emph{the physical} \Phifl, and a $\Delta Y$-only loss does not ask it to. The
  predicted fluxes could be arbitrarily wrong per column while $\Delta Y$ is
  perfect. This substantially weakens the Target-A case \emph{unless} flux-space
  supervision is used, which \S\ref{sec:costwall} showed is affordable only on
  stratified subsets at $\Tn < 5$.
\item \textbf{It interacts with the mask.} Masking a column changes
  $\hat\netflux$ and therefore changes $\nu\hat\netflux$ \emph{only} through the
  component outside $\operatorname{null}(\nu)$. A mask trained under a $\Delta Y$
  loss is being optimised on a projection of its own effect.
\item \textbf{It is a rollout hazard, and a mild one.} Drift within
  $\operatorname{null}(\nu)$ does not affect $\Delta Y$ at all, so it cannot
  destabilise the trajectory directly --- but it does mean the flux diagnostics
  one would use to \emph{understand} an instability are unreliable.
\end{enumerate}

The cheap mitigations are worth naming: regularise $\hat\netflux$ toward small
norm in $\operatorname{null}(\nu)$ (a fixed linear projection, exactly the
machinery of \S\ref{sec:projector} applied to the flux space rather than the
species space); or supervise on \Phifl{} for the subset where it is affordable
and on $\Delta Y$ elsewhere; or report the null-space component of
$\hat\netflux$ as a diagnostic so the drift is at least visible.

Open (register: \reg{R-03}) --- it bears directly on the Target A vs B decision,
which \S\ref{sec:equivalence} showed is not about expressivity. This is a
second, independent cost of Target A, and \S\ref{sec:equivalence}'s trade table
now carries it.

% =====================================================================
\subsection{\texorpdfstring{$\partial R/\partial Y$}{dR/dY} by the product rule}
\label{sec:productrule}

The dominant Jacobian term comes from the reactant product. With
$R_j = \mathrm{base}_j \cdot \prod_{k} Y_{\mathrm{idx}[j,k]}$
($\mathrm{base}_j$ collecting $\lambda$, $\rho$ powers, prefactor):
%
\begin{equation}
D_{ji} \;=\; \frac{\partial R_j}{\partial Y_i}
\;=\; \mathrm{base}_j \sum_{k\,:\,\mathrm{idx}[j,k]=i}\
      \prod_{k'\neq k} Y_{\mathrm{idx}[j,k']}.
\label{eq:dRdY}
\end{equation}

The code implements exactly the sum over slots:

\begin{srclisting}
for k in range(K):
    pe = base.copy()
    for k2 in range(K):
        if k2 != k:
            pe *= Ypad[cn.reactant_idx[:, k2]]
    ...  # accumulate (row=j, col=idx[j,k], value=pe[j])
\end{srclisting}

Two properties this construction has and the obvious alternative does not:

\begin{itemize}
\item \textbf{Repeated reactants are handled by the slot sum.}
  Triple-$\alpha$ has three slots all holding \nuc{He}{4}; the sum of three
  products-of-the-other-two gives $3Y^2$, which is $\dd(Y^3)/\dd Y$. Correct by
  construction, no special case.
\item \textbf{No division.} The tempting shortcut is \code{R\_j / Y\_i}, which is
  faster and wrong: $Y_i = 0$ happens routinely (species below the floor, or
  exactly zero at $t = 0$), giving 0/0. Building the product with the
  differentiated factor \emph{removed} is exact everywhere.
\end{itemize}

The sparse assembly \code{D.tocsr()} then feeds the block Jacobian
%
\begin{equation}
J_{\mathrm{aug}} = \begin{pmatrix}\nu D & 0\\ D & 0\end{pmatrix},
\label{eq:jaug}
\end{equation}
%
whose (2,\,2) block is zero because $R$ does not depend on \Phifl.

\paragraph{A structural observation worth making}\derivedhere: the BDF Newton
matrix is $I - h\beta J_{\rm aug}$, i.e.
%
\begin{equation}
M \;=\; \begin{pmatrix} I_n - h\beta\,\nu D & 0\\
                        -h\beta\,D & I_r\end{pmatrix},
\label{eq:newtonmatrix}
\end{equation}
%
which is \textbf{block lower triangular with an identity block}. Solving
$Mx = b$ is therefore exactly: solve the $n\times n$ system
$(I - h\beta\nu D)x_Y = b_Y$, then $x_\Phifl = b_\Phifl + h\beta D x_Y$ --- no
$r\times r$ factorisation is needed at all. scipy's \code{BDF} does not know
this; it hands the whole $(n{+}r)\times(n{+}r)$ sparse matrix to \code{splu}.
Whether that costs anything real depends on the fill-in the ordering produces,
but it is a concrete optimisation candidate for the \S\ref{sec:costwall} cost
wall, and it is \emph{not} the same as the ADR-0006 contingency (which is about
Jacobian \emph{accuracy}). Logged in \reg{R-20}.

\paragraph{The deliberately neglected chain (ADR 0006).}
$\partial\lambda/\partial Y$ through (i) screening (the plasma state depends on
composition) and (ii) the \Ye-weighting and tabular $\rho\Ye$ inputs is
\emph{not} differentiated. The justification is standard and correct: \textbf{an
approximate Jacobian affects only the Newton convergence rate and hence step
efficiency, never the answer} --- the residual that BDF drives to zero uses the
\emph{exact} RHS. Measured: the neglected chain is $\lesssim 1\%$ of entries
(\code{test\_neglected\_screening\_chain\_is\_subdominant}), and the Jacobian is
validated against finite differences unscreened
(\code{test\_jacobian\_vs\_finite\_differences\_unscreened}).

But \S\ref{sec:costwall} records the bill: at $\Tn \ge 5$ the neglected
\textbf{tabular-EC $\rho\Ye$ chain dominates the weak-drift phase}, the
quasi-Newton iteration stops converging well, and the cost explodes. ``Only
affects efficiency'' is true and can still be the thing that makes a computation
infeasible.

% =====================================================================
\subsection{Conservation by construction, inside the solver}
\label{sec:solverconservation}

The returned composition is \textbf{not} the solver's error-controlled $Y$. It
is

\begin{srclisting}
Y = Y0 + cn.stoich.nu @ Phi
\end{srclisting}

so that:

\begin{itemize}
\item $\Delta Y = \nu\Phifl$ is an \textbf{identity}, not an approximation;
\item baryon and charge conservation hold because $C\nut = 0$
  (\S\ref{sec:targetatheorem}), with the same integer-exactness argument as
  \S\ref{sec:exactzero};
\item the solver's own $Y$ is exposed as a \emph{diagnostic},
  \code{identity\_resid} $= \max|Y_{\rm solver} - Y|$.
\end{itemize}

Measured: \textbf{$2.1\times10^{-16}$ molar} --- machine precision.
\resultsref{2026-07-12}

\begin{invariant}
This is a small design decision with a big property, and it is worth naming the
pattern: \textbf{when two routes to a quantity exist and one is exactly
conserving, return that one and report the other as a residual.} You get
exactness \emph{and} a free consistency check, and the check is sensitive to
anything that would break the correspondence (a wrong $\nu$ column, a solver
bug, an inconsistent RHS).
\end{invariant}

\code{test\_success\_and\_identity\_by\_construction},
\code{test\_solver\_consistency\_residual}, and
\code{test\_conservation\_by\_construction} split those three claims apart.

% =====================================================================
\subsection{Error control for \texorpdfstring{\Phifl}{Phi}}
\label{sec:errorcontrol}

\begin{srclisting}
rtol = 1e-8, atol_y = 1e-18, atol_phi = 1e-18
atol = np.concatenate([np.full(n_s, atol_y), np.full(n_rxn, atol_phi)])
\end{srclisting}

Per-component absolute tolerances, which scipy supports and which matters here.
Tier~0 \S V.6's point applies with force: \textbf{atol matters more than usual}
when the dynamic range is 15+ decades. rtol alone would demand 8 significant
digits on abundances at $10^{-30}$ --- unachievable and pointless.
$\mathrm{atol} = 10^{-18}$ says ``below this, don't care'', and $10^{-18}$ molar
is ${\sim}5\times10^{-17}$ in mass fraction for a mid-network $A \approx 50$ ---
\textbf{about one and a half decades} below the $10^{-15}$ label floor (Tier~0
\S III.6). (Compare the two in the \emph{same} units before quoting a margin:
the three-decade figure is atol in molar against the floor in mass fraction,
which is not a comparison.) One and a half decades is still the right call ---
the floor is chosen relative to what the labels can express, with enough room
that the integrator is never the binding constraint on a quantity the labels
could resolve.

The consequence for cost: the error test runs over \textbf{all} $n + r$
components. For \msn{151} that is 1669 error tests per step, of which 1518 are
\Phifl{} components whose magnitudes span the full range of reaction rates. A
single fast-but-irrelevant column near its atol threshold can throttle the step.
This is the second structural contributor to \S\ref{sec:costwall}, alongside the
Jacobian one.

% =====================================================================
\subsection{The energy identity is a \emph{data} check, not a dynamics check}
\label{sec:energyidentity}

\begin{warn}
\warnglyph{}~\code{CLAUDE.md} invariant \#5 requires the flux route and the
composition route to agree on \enuc{} to $\le 1\%$. It is worth seeing what that
actually tests, because it is not what it looks like.
\end{warn}

The $Q$-value of reaction $j$ is the mass-excess balance:
%
\begin{equation}
Q_j \;=\; -\sum_i \nu_{ij}\,\Delta_i ,
\label{eq:qvalue}
\end{equation}
%
with $\Delta_i$ the mass excess of species $i$ (reactants carry $\nu < 0$, so
this is $\sum_{\rm reactants}\Delta - \sum_{\rm products}\Delta$, as it should
be). Then
%
\begin{equation}
\underbrace{\sum_j Q_j\Phifl_j}_{\text{flux route}}
\;=\; -\sum_j\sum_i \nu_{ij}\Delta_i\Phifl_j
\;=\; -\sum_i \Delta_i(\nu\Phifl)_i
\;=\; \underbrace{-\sum_i \Delta_i\,\Delta Y_i}_{\text{composition route}} .
\label{eq:energycollapse}
\end{equation}

\textbf{The two routes are algebraically identical, given
$\Delta Y = \nu\Phifl$} --- which \S\ref{sec:solverconservation} makes true by
construction. So the residual \textbf{cannot} be measuring integration error.
What it measures is the inconsistency between two nuclear-data sources: the
$Q$-values attached to the rate objects (\code{Rate.Q}, ultimately
REACLIB\,/\,tabular metadata) and the mass excesses attached to the nuclide
objects (\code{Nucleus}). If those were derived from the same mass table to full
precision, the residual would be at machine epsilon.

Measured: median $6\times10^{-6}\ldots2.5\times10^{-5}$ rel, max
$8.7\times10^{-4}$, \textbf{fraction $\le 1\% = 1.000$ in every measured
stratum, both networks, including the whole 3--4\,GK band}.
\resultsref{2026-07-12} Gate passed with $\ge 3$ orders of margin --- and now
you know what the margin is margin \emph{on}: nuclear-data self-consistency at
the $10^{-5}$ level, which is a perfectly reasonable thing to have measured and
a perfectly unreasonable thing to describe as validating the integrator.

\textbf{And the inconsistency is now measured directly.}
\S\ref{sec:wegscheider} runs the test this argument implies --- project $Q$ onto
$\mathrm{rowspace}(\nu)$ and look at the residual --- and finds that the
strong-sector $Q$-values are \textbf{not} reproducible by any single mass table
(rms 0.55\,keV, max 3.5\,keV per column, 331/561 and 972/1345 columns above
0.1\,keV), while the stored $Q$'s differ from pynucastro's own nuclide masses by
up to 17--31\,keV. On $Q \sim 3$--11\,MeV that is 0.2--0.6\% on the worst
channels, which is exactly the shape of the residual distribution above: tiny in
the median, $10^{-3}$-ish in the tail. The energy identity's residual is no
longer an unexplained number.

This does not devalue the check; it relocates it. \textbf{Ask of any green test
what would still be wrong if it passed} (Tier~1 \S VIII.D): here, an integration
error would still be wrong, and this test would not see it. That is what
\code{identity\_resid}, the Radau cross-check, and the label handshake are for.
That relocation is itself a register item --- the gate is real but is not
measuring what its name suggests: \reg{R-23}.

Two known gaps in the identity's coverage, both already on the register:

\begin{itemize}
\item \textbf{Neutrino losses.} Weak reactions emit neutrinos that free-stream
  out (Tier~0 \S0.4.2), so the energy \emph{deposited} is less than $Q$. Whether
  the tabular weak $Q$-values used here are already net-of-neutrino is a
  per-source convention --- Tier~1 \S VIII.C.3's dropped-NU-column item.
\item \textbf{Thermal $Q(T)$.} \code{CLAUDE.md} specifies
  $\enuc = \sum_j Q_j(T)\Phifl_j$ with a \emph{temperature-dependent} $Q$; the
  implementation uses the constant $Q$. Tier~1 \S VIII.E.2 flags it as
  specified-but-unimplemented, and the current gate cannot detect the difference
  precisely \emph{because} of the algebraic collapse above.
\end{itemize}

% ---------------------------------------------------------------------
\subsection{Wegscheider --- the rate set is not globally consistent, measured}
\label{sec:wegscheider}

\begin{warn}
\warnglyph{}~\S\ref{sec:energyidentity} closed by saying the energy residual
measures nuclear-data self-consistency rather than dynamics. This subsection
measures it --- and the measurement has consequences well beyond the energy
gate.
\end{warn}

\subsubsection{The condition}
\label{sec:wegscheider-condition}

Detailed balance constrains rate constants \emph{globally}, not just pairwise.
Around any cycle in the reaction graph the free-energy changes must sum to zero,
so the equilibrium constants must multiply to one \citep{wegscheider1901};
generalised linear-algebraic form --- $\ln K$ must be orthogonal to the null
space of the stoichiometric matrix --- \citep{schuster1989}. In this project's
linear algebra the statement is startlingly simple.

$Q$-values are mass-excess differences: $Q_j = -\sum_i \nu_{ij}\Delta_i$
(\S\ref{sec:energyidentity}). Therefore for \textbf{any} $c$ in the right null
space of $\nu$ --- i.e. any cycle ---
%
\begin{equation}
\sum_j c_j Q_j \;=\; -\sum_j c_j\sum_i \nu_{ij}\Delta_i
\;=\; -\sum_i \Delta_i(\nu c)_i \;=\; 0 .
\label{eq:wegscheider}
\end{equation}

\begin{invariant}
\textbf{Wegscheider condition, linear-algebraic form: $Q$ must lie in the row
space of $\nu$.} Equivalently, $Q \perp \operatorname{null}(\nu)$. If it does
not, no assignment of masses to nuclides reproduces the tabulated $Q$-values,
the network has no globally consistent equilibrium, and \kap{} acquires an
irreducible floor.
\end{invariant}

This is a \emph{free} consistency test on \code{Stoich.Q} --- no rates, no
temperatures, no states. Nothing in the repo performs it.

\subsubsection{The measurement: how far \texorpdfstring{$Q$}{Q} is from the row
               space}
\label{sec:wegscheider-measurement}

Project $Q$ onto $\mathrm{rowspace}(\nu)$ via the SVD and take the residual.
\derivedhere:

\begin{srclisting}
d = np.load("data/stoich/nu_mesa80.npz"); nu, Q = d["nu"], d["Q"]
w = d["weak_mask"].astype(bool)
U, S, Vt = np.linalg.svd(nu[:, ~w], full_matrices=False)
Vr = Vt[:(S > 1e-12 * S[0]).sum()]              # orthonormal basis of rowspace
r  = Q[~w] - Vr.T @ (Vr @ Q[~w])                # the Wegscheider residual
\end{srclisting}

\begin{center}
\begin{tabularx}{\textwidth}{@{}Y L{3.0cm} L{3.0cm}@{}}
\toprule
 & \textbf{\msn{80}} & \textbf{\msn{151}} \\
\midrule
$\|$resid$\|$, \textbf{weak columns only} & $1.8\times10^{-14}$\,MeV &
  $4.4\times10^{-14}$\,MeV \\
$\|$resid$\|$, \textbf{strong\,/\,EM columns} & $1.31\times10^{-2}$\,MeV &
  $1.98\times10^{-2}$\,MeV \\
rms per strong column & $5.53\times10^{-4}$\,MeV &
  $5.39\times10^{-4}$\,MeV \\
max per strong column & $3.47\times10^{-3}$\,MeV &
  $3.57\times10^{-3}$\,MeV \\
strong columns with $|$resid$| > 10^{-4}$\,MeV & \textbf{331\,/\,561} &
  \textbf{972\,/\,1345} \\
$Q$ quantised to $10^{-5}$\,MeV & 94.1\% & 92.1\% \\
\bottomrule
\end{tabularx}
\end{center}

And separately, comparing the stored $Q$ against the \emph{nuclide mass table}
directly ($Q_{\rm pred} = -m_u\cdot\nu\Tr A_{\rm nuc}$):

\begin{center}
\begin{tabularx}{\textwidth}{@{}Y L{3.0cm} L{3.0cm}@{}}
\toprule
 & \textbf{\msn{80}} & \textbf{\msn{151}} \\
\midrule
max $|Q_{\rm stored} - Q_{\rm masstable}|$, strong & \textbf{17.2\,keV} &
  \textbf{30.6\,keV} \\
max $|Q_{\rm stored} - Q_{\rm masstable}|$, weak & 1.1\,keV & 1.1\,keV \\
\bottomrule
\end{tabularx}
\end{center}

\subsubsection{Reading it}
\label{sec:wegscheider-reading}

\textbf{Three separate facts, and they are not the same fact.}

\begin{enumerate}
\item \textbf{The weak sector is exactly consistent} ($10^{-14} =$ machine
  precision). Its $Q$-values come from a single coherent table. Good news, and a
  control: it shows the test is not measuring an artefact of the method.

\item \textbf{The strong sector is not consistent with any mass table
  whatsoever.} The rowspace projection \emph{finds the best-fit mass vector};
  the residual is what survives that fit. rms 0.55\,keV, max 3.5\,keV, and it is
  not round-off: the tabulated $Q$'s are quantised at $10^{-5}$\,MeV $= 10$\,eV,
  so the residual is ${\sim}50$--350 quanta. The top offenders are unmistakable
  ---

\begin{srclisting}
C12_Ne20_to_p_P31      Q =  10.1040   resid = -3.47e-3 MeV
C12_C12_to_He4_Ne20    Q =   4.6210   resid = +2.97e-3
He4_Ne21_to_Mg25       Q =   9.8820   resid = -2.84e-3
\end{srclisting}

  --- REACLIB entries whose $Q$ is quoted to only 3--4 decimals, i.e. older or
  experimentally-derived $Q$-values carried alongside $Q$'s computed from a
  modern mass evaluation. (Forward and reverse of a pair show equal and opposite
  residuals, which is a useful internal check on the method: they share $|Q|$
  and have opposite $\nu$.)

\item \textbf{The stored $Q$'s differ from pynucastro's own nuclide masses by up
  to 31\,keV.} pynucastro's \code{Rate.\_set\_q()} computes $Q$ from
  \code{Nucleus.A\_nuc}, but \code{ReacLibRate} overrides it with the library's
  tabulated $Q$. So the export carries REACLIB's $Q$, and the \emph{mass
  excesses} used elsewhere (NSE in \S\ref{sec:saha}, the composition-route
  energy in \S\ref{sec:energyidentity}) are a different data source.
\end{enumerate}

\subsubsection{Why this matters, quantitatively}
\label{sec:wegscheider-why}

\paragraph{(a) It explains \S\ref{sec:energyidentity}'s residual, and retires it
as a mystery.} The energy identity compares $\sum_j Q_j\Phifl_j$ (REACLIB $Q$)
against $-\sum_i \Delta_i\Delta Y_i$ (nuclide masses). Since the two routes
collapse algebraically given $\Delta Y = \nu\Phifl$, the measured residual
\emph{is} fact (3): 17--31\,keV on $Q \sim 3$--11\,MeV is 0.2--0.6\% on the
worst channels, and the measured energy residual was median
$6\times10^{-6}\ldots2.5\times10^{-5}$ with max $8.7\times10^{-4}$ --- exactly
what you expect when the bad channels are rare and low-flux.
\S\ref{sec:energyidentity} said ``the residual measures nuclear-data
self-consistency''; this subsection measures it.

\paragraph{(b) It predicts a \kap{} floor.} A $Q$ error $\delta Q$ shifts
$\fm/\fp \propto \exp(-Q/kT)$ by a relative $\varepsilon = \delta Q/kT$, and
\S\ref{sec:kappa-db} turns that into $\kap \approx \varepsilon/2$. At
$\Tn = 5$, $kT = 0.431$\,MeV:
%
\begin{equation}
\kap_{\mathrm{floor}} \;\approx\; \frac{\delta Q}{2kT}
\;=\; \begin{cases}
  6.4\times10^{-4} & \delta Q = 0.55\ \text{keV (rms)}\\
  4.0\times10^{-3} & \delta Q = 3.5\ \text{keV (max)}\end{cases}
\label{eq:kappafloor}
\end{equation}
\derivedhere

Now compare with \resultsref{2026-07-10}: \textbf{\kap{} at NSE for stock MESA
r23.05.1 has median $3.6\times10^{-3}$\,/\,$3.3\times10^{-3}$}, recorded there
as attributable to ``own pf-interpolation $+$ winvn provenance'' --- i.e.
attributed loosely and not quantified. The predicted range brackets the measured
value.

\begin{result}
\textbf{This is a hypothesis with a decisive test}, and the test is cheap: run
the rowspace projection on MESA's \emph{own} $Q$ set against MESA's winvn
masses. If the residual is ${\sim}3$--7\,keV, the stock-MESA \kap{} floor is
explained, and the corollary follows immediately --- \textbf{that floor is
irreducible without fixing the mass data, so no \kap{} threshold below
${\sim}10^{-2}$ is measuring physics in MESA-derived quantities.} Given that
\S\ref{sec:negresults}'s \kap-balance census reports ``strong pairs with
$\kap < 10^{-3}$: $\le 0.4\%$ of carrying pairs'', knowing whether $10^{-3}$ is
above or below the data floor is not optional.
\end{result}

Note this is \emph{not} a problem with the project's own engine:
\code{DerivedRate} constructs reverses from the forward using consistent nuclide
data, which is why the engine's \kap{} at NSE collapses to
$2.6\times10^{-12}$ \resultsref{2026-07-10}. The inconsistency does not vanish
--- it is \emph{relocated} into the forward-rate-vs-$Q$ mismatch, where nothing
currently looks for it.

\paragraph{(c) It is the general form of a bug the project already found.}
\resultsref{2026-07-10} records, for MESA 24.08.1, ``a subset of
$(n,\alpha)$\,/\,$(p,\alpha)$ pairs $\kap \sim 0.75$ that are clean in stock ---
consistent with its newer REACLIB snapshot carrying independently-fitted
non-DB-linked pair members; open observation''. That is a Wegscheider violation,
unnamed. The cycle test above is the instrument that would have found it
structurally rather than by noticing an odd \kap{} tail, and would find the next
one.

Open (register: \reg{R-06}). A two-line test that retires one open observation,
quantifies \S\ref{sec:energyidentity}, and may put a floor under a threshold the
kill-test depends on.

% =====================================================================
\subsection{The cost wall, diagnosed}
\label{sec:costwall}

\code{scripts/step6\_integrate\_check.py} measured feasibility on stratified
Sobol states (36\,/\,24, one integration to $10^{2}$\,s serving all nine
measured dts, 11 workers, 4800\,s deadline; censored states are lower bounds).
\resultsref{2026-07-12}:

\begin{center}
\begin{tabularx}{\textwidth}{@{}Y L{3.4cm} L{3.4cm}@{}}
\toprule
 & \textbf{\msn{80}} & \textbf{\msn{151}} \\
\midrule
median wall per state & $\ge 383$\,s & $\ge 1519$\,s \\
censored at 4800\,s & 17/36 & 10/24 \\
\ldots and every censored state is & $\Tn \ge 5$ & $\Tn \ge 5$ \\
throughput & $\le 9.4$ states/h/core & $\le 2.4$ \\
\textbf{full-corpus cost} & \textbf{$\ge$ 110,647 core-h} &
  \textbf{$\ge$ 439,302 core-h} \\
stratified subset ($10^{3}$--$10^{4}$ states, $\Tn < 5$) &
  $\approx 4$--40 core-days & feasible \\
\bottomrule
\end{tabularx}
\end{center}

\textbf{Full-corpus local \Phifl{} supervision is infeasible; stratified subsets
are not.}

And the diagnosis is specific rather than ``it's slow'': the deliberately
neglected tabular-EC $\rho\Ye$ Jacobian chain (\S\ref{sec:productrule})
dominates the weak-drift phase at $\Tn \ge 5$. The physical picture is clean ---
once the strong sector reaches its attractor, all that is left is slow
weak-driven \Ye{} drift through the tabular EC rates, whose $\lambda$ depends on
$\rho\Ye$, i.e. on $Y$, through a chain the Jacobian does not see. The
quasi-Newton iteration then converges poorly exactly in the phase that takes the
most wall time. The ADR-0006 contingency (add the \Ye-chain term) is the named
unlock.

\textbf{Label agreement, where it could be measured:}

\begin{center}
\begin{tabularx}{\textwidth}{@{}L{4.6cm} Y@{}}
\toprule
\textbf{stratum} &
\textbf{median species-fraction in the handshake band} \\
\midrule
QSE window $[3.3, 5.0)$ & 0.74--0.99 (best 0.99 at $[4.0, 5.0)$,
  $\dd t = 10^{-1}$\,s) \\
cold strata & 0.62--0.90 \\
$\Tn \ge 5$ & unmeasurable within deadline \textbf{and} label-pathological \\
\bottomrule
\end{tabularx}
\end{center}

So: \textbf{local \Phifl-label generation is proven for $\Tn < 5$.} The
$\Tn \ge 5$ gap is a double blocker --- cost \emph{and} the fact that the
shipped labels there carry the Appendix-B bug (Tier~3, S13), so agreement would
be uninterpretable anyway.

\begin{repopointer}
\textbf{\textcircled{4} SEE:} notebook \code{11-integrator-and-handoff.ipynb}
--- Fig 1 conservation is exact by construction (not by tolerance), Fig 2 the
cost wall, which is why the node is marked PARTIAL. The notebook deliberately
gates the \code{compile\_network} cells; anything touching it is multi-minute.
\end{repopointer}

% =====================================================================
\subsection{Code and oracles}
\label{sec:codes10}

\textbf{Read:} \code{fluxes/integrate.py} (265 lines). \code{StatePoint} first
--- it hoists everything $Y$-independent for one $(T, \rho)$: the ReacLib set
sums, the pf corrections, $\mathrm{prefactor}\cdot\rho^{\code{dens\_exp}}$. Then
\code{\_lam} (what must be re-evaluated per call: screening, tabular
$\rho\Ye$), \code{gross\_rates}, \code{rate\_jacobian}, \code{integrate\_state}.

The hoisting is the reason this is usable at all: at constant $(T, \rho)$ the
expensive part of Tier~1's machinery is evaluated \textbf{once}, and the RHS is
reduced to a screening call, a tabular interpolation, and a product loop.

\textbf{Oracles} (\code{tests/test\_integrate.py}):

\begin{center}
\begin{tabularx}{\textwidth}{@{}L{6.4cm} Y@{}}
\toprule
\textbf{test} & \textbf{pins} \\
\midrule
\code{test\_rhs\_matches\_engine} &
  the integrator's $R$ \textbf{is} the engine's $R$ ($\le 10^{-12}$ rel) --- no
  second implementation of the physics \\
\code{test\_jacobian\_vs\_finite\_differences\_unscreened} &
  \S\ref{sec:productrule}'s product rule \\
\code{test\_neglected\_screening\_chain\_is\_subdominant} &
  the ADR-0006 approximation, quantified \\
\code{test\_success\_and\_identity\_by\_construction} &
  $\Delta Y = \nu\Phifl$ exactly \\
\code{test\_solver\_consistency\_residual} &
  \code{identity\_resid} at machine precision \\
\code{test\_conservation\_by\_construction} &
  $C\nu = 0$ propagated through the solve \\
\code{test\_something\_burned} &
  the anti-vacuity test --- the integration actually did something \\
\code{test\_energy\_identity} & \S\ref{sec:energyidentity} \\
\code{test\_t\_eval\_multiple\_dts} &
  one integration serves all nine label dts \\
\code{test\_vs\_shipped\_labels} & the external comparison \\
\bottomrule
\end{tabularx}
\end{center}

\code{test\_something\_burned} deserves a note: it is the guard against a test
suite that passes on a null integration. Every identity above holds trivially
for $\Phifl = 0$. Without it, a bug that made the RHS return zeros would pass
eight of the nine tests.

% =====================================================================
\subsection{Self-check}
\label{sec:selfcheck-s10}

\begin{selfcheck}
\item Write $J = \nu D$ and explain why the stiffness ratio here exceeds
  $10^{14}$, citing the measurement.
\item Backward Euler and the trapezoidal rule are both A-stable. Why is only one
  of them acceptable here?
\item Why is \Phifl{} carried as a state variable rather than reconstructed from
  $\Delta Y$ afterwards? Give the dimension of the ambiguity.
\item Why does \code{rate\_jacobian} build products with the differentiated
  factor removed instead of dividing?
\item \code{integrate\_state} returns $Y_0 + \nu\Delta\Phifl$ rather than the
  solver's $Y$. What is gained, and what does \code{identity\_resid} then
  measure?
\item Show that $\sum_j Q_j\Phifl_j = -\sum_i \Delta_i \Delta Y_i$ identically.
  What, therefore, does the 1\% energy gate actually test?
\item The augmented Newton matrix is block lower triangular. What does that
  imply about the cost of carrying \Phifl, and does the current implementation
  exploit it?
\item Full-corpus \Phifl{} labels are infeasible. What is the measured cause,
  and what is the named unlock?
\end{selfcheck}
